% ID: SEL-D0056-OA-E02
% ESTADO: cotejada. Segunda pasada visual y borrador compilado revisados.
% Fuente: Selectividad/sel_2011_matematicas_aplicadas/123-EXAMEN 1.pdf
% SHA-256: 69f8b84ed22f2b49074eafb95dbdce8a98114a4a2f02989b92f3f7add9355155
% Página: 1; opción A; ejercicio original 2.

Tras un test realizado a un nuevo modelo de automóvil, se ha observado que el consumo de gasolina, $c(x)$, expresado en litros, viene dado por la función
\[
c(x)=7.5-0.05x+0.00025x^2,
\]
siendo $x$ la velocidad en km/h y $25\leq x\leq175$.
\begin{enumerate}[label=\alph*)]
\item \textbf{(0.5 puntos)} Determine el consumo de gasolina a las velocidades de 50 km/h y 150 km/h.
\item \textbf{(1 punto)} Estudie el crecimiento y decrecimiento de la función $c(x)$.
\item \textbf{(1 punto)} \textquestiondown{}A qué velocidades de ese intervalo se obtiene el mínimo consumo y el máximo consumo y cuáles son éstos?
\end{enumerate}
